\section{Changing orbit weights and locating the trace}\label{prof:orbits}

This section changes from the counting measure to Haar measure on an
orbit that fixes the three distinguished root directions. It then places
the Gaussian trace within one length standard deviation, using three
length densities rather than a direct trace expansion.

\subsection{The determinant weight}

Fix the root $y$ and the conditional matrix data modulo proper rotations
that fix $\mathcal P$ pointwise. For a proper rotation $O$ of
$\mathcal P^\perp$, put $R=I_{\mathcal P}\oplus O$. The orbit action is
\[
 J\longmapsto RJR^{\mathsf T},\qquad
 h=\Phi_J(y)\longmapsto Rh.
\]
It fixes $y,m,\mathbf1$ and preserves the root equation. The original
spectral restriction is constant on the orbit. Before the determinant
weight, disintegrate the root integral as $d\pi(\omega)\,d\mathrm{Haar}(O)$,
where $\omega$ denotes the base data and $d\pi$ includes the scalar
factors and, when imposed, the invariant spectral restriction.
The conditional matrix law is Haar in $O$ because all its fixed means
and axes lie in $\mathcal P$.

Let $\mathcal D_\omega$ be the set of orientations satisfying all
precision, trace, count, weighted, and projection diagnostics of
Proposition~\ref{prof:conditioned-precision} with the chosen tolerances.
This is stronger than mere Jacobian regularity. Call an orbit suitable
when $\mathrm{Haar}(\mathcal D_\omega^c)\le e^{-B_0k}$ for a large
fixed $B_0$. The integrated matrix exclusion and Markov's inequality
remove unsuitable base data at any desired smaller rate, by starting
with a larger raw rate. Lemma~\ref{prof:determinant-square} and
Cauchy--Schwarz give the same conclusion after determinant weighting.

On $\mathcal D_\omega$, $K_T\succeq cq^2I$. Extend scalar $\log t$
linearly below $cq^2$, and let $F(O)$ be the trace of this extended
function of $K_T$. It equals $\log\det K_T$ on $\mathcal D_\omega$.
Concavity of the extension and the trace-norm bound on the second
rotation derivative of $J$ give the upper Taylor bound
\[
 F(O')\le F(O)+\langle\nabla F(O),\dot O\rangle
                    +Cq^{-2}\operatorname{dist}(O,O')^2
\]
along a minimizing geodesic, with the linear term evaluated on its
initial tangent. On diagnostic orientations,
$\|\nabla F\|\le\|[J,K_T^{-1}]\|_{\rm HS}=o(\sqrt n)$.
For the diagonally truncated precision, the commutator is bounded by
$Cq^{1-\eta}\sqrt{n/q}$. Restoring its positive diagonal changes the
inverse by a matrix of the discarded rank; the weighted inverse and
operator bounds give commutator cost at most
$Cq^{-1-\eta}\sqrt{\operatorname{rank}}$. Since that rank is
$O(nq^3)$ and $\eta$ is small, both costs are $o(\sqrt n)$.

Choose a median $m_F$ of $F$ on $\mathcal D_\omega$. Its high and low
anchor sets each have a fixed positive Haar mass. Distance to either
set has a sub-Gaussian tail at scale $n^{-1/2}$. The upper Taylor bound
from a low anchor bounds $F-m_F$ above; applying it from a diagnostic
orientation to a high anchor bounds $F-m_F$ below. At distance
$C_B\sqrt{k/n}$, the two errors are
\[
 O_B\left(\sqrt k+\frac{k}{nq^2}\right)=o(k),
\]
and the complementary Haar mass is $e^{-Bk}$. At distance $O(n^{-1/2})$
the errors are bounded. Integrating the same distance tail shows that
$\E_{\rm Haar}e^{t(F-m_F)}$ is bounded for every fixed $t>0$.
This also controls the weighted contribution of exceptional diagnostic
orientations.

Write $w=\mathcal W(y,J)$ and
\[
 w_0=e^{-nb^2/2}\det(V)e^{m_F}.
\]
On $\mathcal D_\omega$, $w=e^{-nb^2/2}\det(V)e^{F}$, so the preceding
bounds compare $w$ with the constant orbit weight $w_0$. A fixed
positive Haar fraction has $w\ge c w_0$; the total root-count bound
therefore gives $\int w_0\,d\pi\le n^C$ on suitable base data.
For every $0\le p(O)\le1$, every small fixed $\xi>0$, and every
needed fixed $B$, the comparison yields
\[
 \int w_0\E_{\rm Haar}p\,d\pi
 \le e^{\xi k}\int\E_{\rm Haar}(wp)\,d\pi
       +e^{-Bk}\int w_0\,d\pi.
\]
Conversely,
\[
 \int\E_{\rm Haar}(wp)\,d\pi
 \le e^{\xi k}\int w_0\E_{\rm Haar}p\,d\pi
       +e^{-Bk}\int w_0\,d\pi+\mathcal E_B,
\]
where $\mathcal E_B\le n^Ce^{-Bk}$ is the globally excluded weighted
mass of non-diagnostic orientations, after the raw matrix rate is
chosen. The comparison is integrated
over $\omega$; it does not require a lower weight on each spectral
fiber.

These inequalities permit further Markov exclusions over base data.
When the input is a field probability with a planted spin, the degree
identity in Corollary~\ref{prof:unique-root} charges excess multiplicity
to twice the negative-root count. Its integrated exclusion rate is
already available. Thus the same comparisons apply to those posterior
tests after this multiplicity error is included.

\subsection{Median trace control}

We need one further estimate on most orientations of each suitable
orbit:
\begin{equation}\label{prof:eq12}
 |\operatorname{Tr}K^{-1}-nb|\le C\sqrt{n/q}.
\end{equation}
The precision estimates give the fluctuation scale of Gaussian length,
but not the location of its mean. We compare its density at the target
and one fluctuation on either side. The lemma below shows that these
three densities locate the trace.

Put $L=\sqrt{n/q}$ and $y_w=nb+wL$ for $w=0,\pm1$; here $y_w$ is a
scalar squared length. On $\mathcal P^\perp$, let $\gamma_V$ be the
normalized law $N(0,V)$ conditioned on zero projection onto
$\mathcal P$, and let $\gamma_b=N(0,bI_{\mathcal P^\perp})$. Denote
their squared-length densities by $\ell_V$ and $\ell_b$. Define on
every orientation
\[
 A_w(O)=\ell_V(y_w)\,
 \E_{\gamma_V}\!\left[
 e^{z^{\mathsf T}(J-bI)z/2}\mid\|z\|^2=y_w\right].
\]
This is an unnormalized weighted length density. On a diagnostic
orientation the full weighted integral is finite; writing its mass as
$Z_V^\perp(O)$ gives
\[
 A_w(O)=Z_V^\perp(O)f_{\widehat C_O}(y_w),\qquad
 \widehat C_O=(K|_{\mathcal P^\perp})^{-1},
\]
where $f_C$ denotes the squared-length density under $N(0,C)$.
The fixed-length definition of $A_w$ remains meaningful even on the
other orientations.

For the scalar prior, the normalized weighted covariance is
\[
 \widehat C_o=\bigl(((b+b^{-1})I-J)|_{\mathcal P^\perp}\bigr)^{-1}
             =(K_o|_{\mathcal P^\perp})^{-1}.
\]
Let $Z_o^\perp$ be its unnormalized weighted mass and
$f_o=f_{\widehat C_o}$. Both are invariant under the remaining
rotation. Haar averaging at each fixed length gives the exact identity
\[
 \overline A_w:=\E_O A_w(O)
 =Z_o^\perp f_o(y_w)\frac{\ell_V(y_w)}{\ell_b(y_w)}.
\]
The ratios $\overline A_{\pm1}/\overline A_0$ are bounded above and
below. For $f_o$, this follows from
$\operatorname{Tr}\widehat C_o=nb+o(L)$, its trace-square scale $L^2$,
and the Fourier local limit supplied by the spread. For the prior
ratio, both traces are $nb+O(1)$, the top eigenvalues are bounded, and
the trace-square discrepancy is $O(nq^2+1)$. The scalar length-tilt
expansion therefore gives
\[
 \log\frac{\ell_V(y_w)/\ell_b(y_w)}
          {\ell_V(y_0)/\ell_b(y_0)}
 =O\!\left(\frac{L^2q^2}{n}+\frac{L^3}{n^2}+o(1)\right)=O(1).
\]
The normalizer $Z_o^\perp$ is the same at all three lengths.

We next show $\E_O(A_w/\overline A_w)^2\le1+o(1)$. Let
$\rho_V(T\mid y_w,y_w)$ and $\rho_b(T\mid y_w,y_w)$ be the overlap
densities of two independent draws from the corresponding unweighted
priors, conditioned on those lengths. Let $\mu_{o,w}$ be the scalar
prior at length $y_w$, weighted by the common angular energy and
normalized. Haar averaging of a pair at fixed overlap gives
\[
 \E_O\left(\frac{A_w}{\overline A_w}\right)^2
 =\E_{\mu_{o,w}^{\otimes2}}
   \frac{\rho_V(T\mid y_w,y_w)}{\rho_b(T\mid y_w,y_w)},
 \qquad T=z_1^{\mathsf T}z_2.
\]
Thus the positive multiplier is a ratio of \emph{prior} densities,
while the expectation and its tails use the \emph{energy-tilted}
reference law.

To estimate the multiplier, tilt each prior length to $y_w$; its
scalar tilt is $O(L/n)$. At zero cross tilt the overlap is symmetric,
and its covariance with either length is zero. Its conditional
curvature is therefore $\operatorname{Tr}\widetilde V^2$, where
$\widetilde V$ is the length-tilted covariance. The two tilted traces
are the same target, and their trace squares differ relatively by
$O(q^2+1/n)$: the variable prior differs from scalar covariance by
normalized Hilbert--Schmidt size $O(q+1/\sqrt n)$, and the scalar tilt
changes eigenvalues by a uniformly Lipschitz function. Analytic
expansion of the three quadratic log moments, with their nondegenerate
Hessian, and Fourier inversion of the central prefactors give, for
$|T|/n$ small,
\[
 \log\frac{\rho_V(T\mid y_w,y_w)}{\rho_b(T\mid y_w,y_w)}
 \le o(1)+C(q^2+1/n)\frac{T^2}{n}+C\frac{T^4}{n^3}.
\]
Odd cross derivatives vanish by symmetry; the prefactor error is
$O(q+|T|/n)$ in the same small-scale convention. Conditioning the
initial projections at zero changes only a fixed rank.

Here is the tail comparison that makes this local bound integrable.
Put $U=T/L$ and $N=nq^3$. The positive quadratic and quartic errors
become
\[
 C\left((q+(nq)^{-1})U^2+\frac{U^4}{nq^2}\right).
\]
Under $\mu_{o,w}^{\otimes2}$, quadratic exponentiation and the two
length-density bounds give a Gaussian tail in $U$ up to
$|U|\le c\sqrt N$. There the quartic error is at most $C c^2qU^2$,
so the Gaussian tail dominates after the upper scale is decreased;
for fixed $U$ the multiplier tends to one. For
$a=|T|/n$ from order $q$ to a small fixed cutoff, the reference
length-deficit estimate of Section~\ref{prof:subsec:replica-integrals}
gives $e^{-cna^3}$. Dividing the positive exponent by $na^3$ leaves
$O(q^2/a+a+1/(na))$, which is small in this range after the cutoff
and upper scale are chosen. These two ranges give uniform integrability
through small overlaps.

For a fixed larger overlap, simultaneous quadratic Chernoff bounds for
the unweighted priors, conditional on both lengths, give at least the
spherical fixed-overlap rate $-\tfrac12\log(1-a^2)$ up to an $o(1)$
loss per coordinate as the upper scale decreases. This rate is strictly
larger than $a^2/2$ away from zero. Near the parallel endpoints, fixed
tilts chosen before $n$ supply the finite strict rate gap needed below.
Exact length densities under these tilts cost at most a polynomial by
Fourier inversion, using damping from typical sites or a bounded-eigenvalue
block; the preceding $O(L)$ prior length deviations cost subexponentially.
The angular pair upper bound of
Section~\ref{prof:subsec:replica-integrals} costs only $a^2/2$ after
normalization: the reference angular integral per draw is at least
$\exp\{n(1/4-o(1))\}$. To verify that lower bound, use the Gaussian
formula on $\mathcal P^\perp$ at precision
$((b+b^{-1})I-J)|_{\mathcal P^\perp}$. Its centered length density
costs only a polynomial, its normalized log determinant is at most
$1/2+o(1)$ by the semicircle bound and fixed-rank interlacing, and
$y_w/n=1+o(1)$. The strict rate gap controls the remaining overlaps.
Thus the displayed second moment is $1+o(1)$.

The first and second moments imply that all three ratios
$A_w/\overline A_w$ are close to one outside as small a fixed Haar
mass as desired. On diagnostic orientations the common factor
$Z_V^\perp(O)$ cancels, so
$f_{\widehat C_O}(y_{\pm1})/f_{\widehat C_O}(y_0)$ is bounded above.
The next lemma now applies. Its operator and trace-square hypotheses
follow from Proposition~\ref{prof:conditioned-precision} after removing
three dimensions. Since $v_i\ge1/2$ outside at most $2nq$ sites, the
coordinate subspace supported on $\{i:v_i\ge1/2\}$ has intersection
with $\mathcal P^\perp$ of dimension at least $n-2nq-3$. The compressed
precision has bounded norm on this intersection, so the min--max principle
gives a fixed positive fraction of covariance eigenvalues bounded below.
Finally, fixed-rank interlacing gives
$|\operatorname{Tr}K^{-1}-\operatorname{Tr}\widehat C_O|
\le Cq^{-2}=o(L)$. The lemma proves \eqref{prof:eq12}.

\begin{lemma}[Three length densities locate a Gaussian trace]
\label{prof:trace-from-three-densities}
Fix positive constants $c_0,C_0,M$. Let $0<q\le q_0$ and
$nq^3\ge\log^{10}n$, and put $L=\sqrt{n/q}$. Let $C$ be a positive
definite matrix of dimension $n-O(1)$ such that
\[
 \|C\|\le C_0q^{-2},\qquad
 c_0L^2\le\operatorname{Tr}C^2\le C_0L^2.
\]
Assume also that at least $c_0n$ eigenvalues of $C$ exceed $c_0$.
Let $f$ be the density of $\|C^{1/2}g\|^2$, with $g$ standard Gaussian,
and let $x_0\in[c_0n,C_0n]$. If
\[
 f(x_0+L)\le M f(x_0),\qquad
 f(x_0-L)\le M f(x_0),
\]
then $|\operatorname{Tr}C-x_0|\le C_1L$, where $C_1$ depends only on
the fixed constants. The assertion holds for all sufficiently large $n$.
\end{lemma}

\begin{proof}
Write $m(\lambda)=\operatorname{Tr}(C^{-1}+\lambda I)^{-1}$ on
$\lambda>-\|C\|^{-1}$ and $C_\lambda=(C^{-1}+\lambda I)^{-1}$.
The function $m$ decreases continuously from infinity to zero, and
\[
 m'(\lambda)=-\operatorname{Tr}C_\lambda^2.
\]
If $|m(0)-x_0|\le L$, there is nothing to prove. Otherwise set
$\varepsilon=\operatorname{sign}(m(0)-x_0)$ and
$x_1=x_0+\varepsilon L$. There is a unique $\lambda$ with
$m(\lambda)=x_1$, and $\varepsilon\lambda>0$.
Exponential tilting of the squared length gives its density $f_\lambda$
under covariance $C_\lambda$ and the identity
\[
 \frac{f(x_1)}{f(x_0)}
 =e^{|\lambda|L/2}\frac{f_\lambda(x_1)}{f_\lambda(x_0)}.
\]
This identity treats the two signs with exactly the same inequality.

First we bound the tilt without assuming a local central limit theorem
near the final pole. This coarse bound will place the covariance in a
region comparable to its original value; sharper density estimates can
then improve the logarithmic tilt bound to a constant. Since
$\operatorname{Tr}C_\lambda=x_1\asymp n$, its standard deviation is
at most $\sqrt2x_1\le Cn$. The mean-density bound of
Lemma~\ref{prof:quadratic-mean-density} gives
$f_\lambda(x_1)\ge c/n$.
If $\lambda>0$, the inequality $m(\lambda)\le n/\lambda$ implies
$\lambda\le C$; if $\lambda<0$, every eigenvalue increases.

In both cases at least $c n$ eigenvalues of $C_\lambda$ are bounded
below by a fixed positive constant. The quadratic Gaussian product
formula therefore gives
\[
 |\mathbb E e^{it\|C_\lambda^{1/2}g\|^2}|
 \le (1+c t^2)^{-cn},
 \qquad \sup_x f_\lambda(x)\le Cn^{-1/2}.
\]
Consequently
$M\ge c n^{-1/2}e^{|\lambda|L/2}$, so
$|\lambda|L\le C\log n$.
This is the promised control near a pole: no local central limit theorem
is being assumed there, and the density loss is only polynomial.

Choose $a>0$ so small that $aC_0<1/2$. Because
$q^2L=\sqrt{nq^3}\ge\log^5 n$, the preceding bound forces
$|\lambda|<a q^2$. Thus $C_\theta$ is comparable to $C$ for every
$\theta$ between zero and $\lambda$, and
$\operatorname{Tr}C_\theta^2\asymp L^2$.
We can now use the spectral spread at the tilted covariance. For these
tilts, the largest eigenvalue divided by the square root of
the trace square is $O((nq^3)^{-1/2})=o(1)$.
To check the Fourier bound explicitly, if $a_i$ are the eigenvalues of
$C_\lambda$ and $\sigma^2=2\sum_i a_i^2$, concavity of
$x\mapsto\log(1+ux)$ on $[0,\max_i a_i^2]$ gives
\[
 \prod_i(1+4t^2a_i^2/\sigma^2)^{-1/4}
 \le\left(1+\frac{4t^2a_{\max}^2}{\sigma^2}\right)
       ^{-\sigma^2/(8a_{\max}^2)}.
\]
Since $a_{\max}^2/\sigma^2\le1/16$ for large $n$, the right-hand side is
bounded by $(1+t^2/4)^{-2}$, an integrable function independent of the
parameters.

Hence its density has supremum at
most $C/L$. Its density at the mean is at least $c/L$ by the same
mean-density lemma. Applying these bounds to $f_\lambda$ in the ratio
identity improves the preceding estimate to $|\lambda|L\le C$.
Finally,
\[
 |m(0)-x_0|
 \le L+|m(0)-m(\lambda)|
 \le L+\int_0^{|\lambda|}C L^2\,du\le C_1L.
\]
\end{proof}
